%%% Hlavní soubor. Zde se definují základní parametry a odkazuje se na ostatní části. %%%

%% Verze pro jednostranný tisk:
% Okraje: levý 40mm, pravý 25mm, horní a dolní 25mm
% (ale pozor, LaTeX si sám přidává 1in)
\documentclass[12pt,a4paper]{report}
\setlength\textwidth{145mm}
\setlength\textheight{247mm}
\setlength\oddsidemargin{15mm}
\setlength\evensidemargin{15mm}
\setlength\topmargin{0mm}
\setlength\headsep{0mm}
\setlength\headheight{0mm}
% \openright zařídí, aby následující text začínal na pravé straně knihy
\let\openright=\clearpage

%% Pokud tiskneme oboustranně:
% \documentclass[12pt,a4paper,twoside,openright]{report}
% \setlength\textwidth{145mm}
% \setlength\textheight{247mm}
% \setlength\oddsidemargin{15mm}
% \setlength\evensidemargin{0mm}
% \setlength\topmargin{0mm}
% \setlength\headsep{0mm}
% \setlength\headheight{0mm}
% \let\openright=\cleardoublepage

%% Použité kódování znaků: obvykle latin2, cp1250 nebo utf8:
%\usepackage[utf8]{inputenc}
%\usepackage[T2A]{fontenc}
%\usepackage[russian,czech,english]{babel}

%\usepackage{fontspec}

%\usepackage{xunicode,fontspec,xltxtra}
%\usepackage[english]{polyglossia}
%\setotherlanguages{russian} % set as "other" so English hyphenation active


\usepackage{libertine}


%% Ostatní balíčky
\usepackage{verbatim}
\usepackage{graphicx}
%\usepackage{amsthm}

%% Balíček hyperref, kterým jdou vyrábět klikací odkazy v PDF,
%% ale hlavně ho používáme k uložení metadat do PDF (včetně obsahu).
%% POZOR, nezapomeňte vyplnit jméno práce a autora.
\usepackage[unicode]{hyperref}   % Musí být za všemi ostatními balíčky
\hypersetup{pdftitle=A Comparison of Methods of Czech-to-Russian Machine Translation}
\hypersetup{pdfauthor=Karel Bílek}

\usepackage{tabularx}

\usepackage{polyglossia}
\setmainlanguage{english}
\setotherlanguage{czech}

\usepackage{float}

\usepackage[
   backend=bibtex8      % if we want unicode 
  ,style=iso-authoryear % or iso-numeric for numeric citation method          
  ,babel=other        % to support multiple languages in bibliography
  ,sortlocale=cs_CZ   % locale of main language, it is for sorting
  %,bibencoding=UTF8   % this is necessary only if bibliography file is in different encoding than main document
]{biblatex}


\bibliography{library}
%%% Drobné úpravy stylu
% Tato makra přesvědčují mírně ošklivým trikem LaTeX, aby hlavičky kapitol
% sázel příčetněji a nevynechával nad nimi spoustu místa. Směle ignorujte.
\makeatletter
\def\@makechapterhead#1{
  {\parindent \z@ \raggedright \normalfont
   \Huge\bfseries \thechapter. #1
   \par\nobreak
   \vskip 20\p@
}}
\def\@makeschapterhead#1{
  {\parindent \z@ \raggedright \normalfont
   \Huge\bfseries #1
   \par\nobreak
   \vskip 20\p@
}}
\makeatother

% Toto makro definuje kapitolu, která není očíslovaná, ale je uvedena v obsahu.
\def\chapwithtoc#1{
\chapter*{#1}
\addcontentsline{toc}{chapter}{#1}
}


\newfloat{tabulka}{tbph}{lop}
\floatname{tabulka}{Table}

\long\def\centergraf#1 {
{
    \begingroup
    \centering
    \par
    #1
    \par
    \endgroup
}

}

\def\uv#1{
 ``#1''
}

\long\def\jednatabulkan#1#2#3#4{
 \begin{tabulka}[h]
     
         \centergraf{
         \footnotesize
    \begin{tabular}{ #2 }
        #3
        
    \end{tabular}
             
             }
     
     \caption{#4} 
     \label{tabulka:#1}
 \end{tabulka}
}



\long\def\jednatabulka#1#2#3#4{
 \begin{tabulka}[h]
     
         \centergraf{
         \footnotesize
    \begin{tabularx}{\textwidth}{ #2 }
        #3
        
    \end{tabularx}
             
             }
     
     \caption{#4} 
     \label{tabulka:#1}
 \end{tabulka}
}

\begin{document}

% Trochu volnější nastavení dělení slov, než je default.
\lefthyphenmin=2
\righthyphenmin=2

%%% Titulní strana práce

\pagestyle{empty}
\begin{center}

\large

Charles University in Prague

\medskip

Faculty of Mathematics and Physics

\vfill

{\bf\Large MASTER THESIS}

\vfill

\centerline{\mbox{\includegraphics[width=60mm]{../img/logo.pdf}}}

\vfill
\vspace{5mm}

{\LARGE Karel Bílek}

\vspace{15mm}

% Název práce přesně podle zadání
{\LARGE\bfseries A Comparison of Methods of Czech-to-Russian Machine Translation}

\vfill

% Název katedry nebo ústavu, kde byla práce oficiálně zadána
% (dle Organizační struktury MFF UK)
%Ústav formální a aplikované lingvistiky
Institute of Formal and Applied Linguistics

\vfill

\begin{tabular}{rl}

Supervisor of the master thesis: & doc. RNDr. Vladislav Kuboň, Ph.D. \\
\noalign{\vspace{2mm}}
Study programme: & Informatika  ????  \\
\noalign{\vspace{2mm}}
Specialization: & Asi něco s lingvistikou \\
\end{tabular}

\vfill

% Zde doplňte rok
Prague 2014

\end{center}

\newpage

%%% Následuje vevázaný list -- kopie podepsaného "Zadání diplomové práce".
%%% Toto zadání NENÍ součástí elektronické verze práce, nescanovat.

%%% Na tomto místě mohou být napsána případná poděkování (vedoucímu práce,
%%% konzultantovi, tomu, kdo zapůjčil software, literaturu apod.)

\openright

\noindent
Dedication. I love everybody. (TODO)

\newpage

%%% Strana s čestným prohlášením k diplomové práci

\vglue 0pt plus 1fill

\noindent
I declare that I carried out this master thesis independently, and only with the cited
sources, literature and other professional sources.

\medskip\noindent
I understand that my work relates to the rights and obligations under the Act No.
121/2000 Coll., the Copyright Act, as amended, in particular the fact that the Charles
University in Prague has the right to conclude a license agreement on the use of this
work as a school work pursuant to Section 60 paragraph 1 of the Copyright Act.

\vspace{10mm}

\hbox{\hbox to 0.5\hsize{%
In ........ date ............
\hss}\hbox to 0.5\hsize{%
signature of the author
\hss}}

\vspace{20mm}
\newpage

%%% Povinná informační strana diplomové práce

\vbox to 0.5\vsize{
\setlength\parindent{0mm}
\setlength\parskip{5mm}

Název práce:
Porovnáni metod česko-ruského automatického překladu

% přesně dle zadání

Autor:
Karel Bílek

Katedra:  % Případně Ústav:
Ústav formální a aplikované lingvistiky
% dle Organizační struktury MFF UK

Vedoucí diplomové práce:
doc. RNDr. Vladislav Kuboň, Ph.D.,
Ústav formální a aplikované lingvistiky
%Jméno a příjmení s tituly, pracoviště
% dle Organizační struktury MFF UK, případně plný název pracoviště mimo MFF UK

Abstrakt:
Abstrakt abstrakt.
% abstrakt v rozsahu 80-200 slov; nejedná se však o opis zadání diplomové práce

Klíčová slova:
% 3 až 5 klíčových slov
Čeština, ruština, strojový překlad

\vss}\nobreak\vbox to 0.49\vsize{
\setlength\parindent{0mm}
\setlength\parskip{5mm}

Title:
% přesný překlad názvu práce v angličtině
{A Comparison of Methods of Czech-to-Russian Machine Translation}

Author:
Karel Bílek
%Jméno a příjmení autora

Department:
Institute of Formal and Applied Linguistics
%Ústav formální a aplikované lingvistiky
%Institute of Formal and Applied Linguistics
%Název katedry či ústavu, kde byla práce oficiálně zadána
% dle Organizační struktury MFF UK v angličtině

Supervisor:
doc. RNDr. Vladislav Kuboň, Ph.D.,
Institute of Formal and Applied Linguistics
%Jméno a příjmení s tituly, pracoviště
% dle Organizační struktury MFF UK, případně plný název pracoviště
% mimo MFF UK v angličtině

Abstract:
Abstrakt abstrakt.
% abstrakt v rozsahu 80-200 slov v angličtině; nejedná se však o překlad
% zadání diplomové práce

Keywords:
Czech, Russian, Machine translation
% 3 až 5 klíčových slov v angličtině

\vss}

\newpage

%%% Strana s automaticky generovaným obsahem diplomové práce. U matematických
%%% prací je přípustné, aby seznam tabulek a zkratek, existují-li, byl umístěn
%%% na začátku práce, místo na jejím konci.

\openright
\pagestyle{plain}
\setcounter{page}{1}
\tableofcontents

%%% Jednotlivé kapitoly práce jsou pro přehlednost uloženy v samostatných souborech
\chapter*{Introduction}
\addcontentsline{toc}{chapter}{Introduction}
Slavic  family of languages consists of around 12 languages, de\-pen\-ding on the exact definition.\footnote{See for example \cite{sussex2011slavic} or \cite{siewierska1998overview} --  the final language count depends on whether Upper and Lower Sorbian are taken as different languages or not; whether Kashubian is a separate language, or just a variety of Polish; whether Serbian and Croatian are treated as separate languages or not}
They are usually divided into three sub-groups - South Slavic, West Slavic and East Slavic. While languages inside each of those groups are mutually intelligible, intelligibility accross those groups is lower; still, the Slavic languages in general have very similar morphology, syntax and vocabulary.

Czech language belongs to the West Slavic group, while Russian belongs to the East Slavic group.  Therefore, they are not fully mutually intelligible -- but, as noted, as Slavic languages, they are in some respects still very similar. 
This creates a great environment for building machine translations tools -- the languages are not too close to make the task meaningless, while we can still take an advantage of the closeness.

For political reasons, significant efforts were directed towards building an automatic translation system between Russian and Czech, mainly during mid-eighties with the RUSLAN system. The efforts were terminated in 1990 in the final phases of development.\footnote{See for example \cite{recycled}, \cite{hajic1987}, \cite{olivaruslan}} 

After the velvet revolution, political dependence on Russian-speaking countries were weakened, while the economical and cultural dependence on English-speaking countries were strengthened. Therefore, the efforts of automatic translation were more focused on English. 

In the process, general direction of machine translation efforts shifted from rule-based systems to more statistical-based systems. This allows us to compare older, rule-based approaches to the more modern statistical approaches.

In this thesis, I am trying to describe the rich history of automatic Czech-to-Russian translation, then experiment with the available tools -- both historical and more recent ones -- on a given set of data, and finally compare the results.

\chapter{Brief history of Czech and Russian}

\section{Proto-Slavic}

The Indo-Europeans occupied European lands from approximately 4000 BC. Their language, Proto-Indo-European (PIE), is not known and can only be linguistically reconstructed. 

In approximately 2000 BC, a Proto-Slavic (PS) language started to emerge. Until approximately 400 AD, the language was fairly uniform and the land occupied by Slavs was broadly coherent, stretched approximately from Oder river (Odra in Polish) to Dnieper (Днепр in Russian).\footnote{See \cite{sussex2011slavic}, pg.19 and further} 

In this part, I am describing some of the general changes from PIE to PS.
I am focusing mainly on syntactical changes -- not because there are no morphological, phonological or lexical changes, but because there are too many to cover here.

If a reader is interested in a further study of a history of Slavic languages and their difference in total, I can only recommend \cite{sussex2011slavic}. \footnote{Unless otherwise noted, information in this chapter are taken from this book.}

\subsection{Cases}
\jednatabulka{korintskym} { |l|X |X | X | }
{
         \hline
genitive &
invocant nomen \textbf{Domini} nostri Iesu Christi
&
   who call on the name \textbf{of} our \textbf{Lord} Jesus Christ         
&
všem, kdo na jakémkoli místě vzývají jméno našeho společného \textbf{Pána} Ježíše Krista
\\

    \hline
ablative
&
gratia vobis et pax a Deo Patre nostro et \textbf{Domino} Iesu Christo
&
Grace and peace to you \textbf{from} God our Father and \textbf{the Lord} Jesus Christ
&
Milost vám a pokoj od Boha, našeho Otce, a od \textbf{Pána} Ježíše Krista \\
    \hline
} {Demonstrating ablative and genitive conflation on 1 Corinthians} 


PIE had seven cases - nominative, accusative, genitive, dative, instrumental, locative, ablative and vocative. 
In PS, ablative and genitive were conflated into just genitive. 

To illustrate, I have added a simple comparison of old latin and Czech. 

Latin retained both ablative and genitive. Latin genitive of \emph{dominus} (lord, master) is \emph{dominī}, latin ablative of the same word is  \emph{dominō}. Both these forms were used in the very beginning of 1~Corinthians in Latina Vulgata (latin version of The Bible).

In the Table~\ref{tabulka:korintskym}, I have shown the translation to both Czech and English. We can see that both cases are inflated in Czech as genitive \uv{(našeho) pána}.\footnote{Bible sources: \cite{latinavulgata}, \cite{bibleniv}, \cite{bible21}. Note that English and Czech translations are not actually translations from Latin, but from more primary sources, but it will suffice for this simple comparison.}

\subsection{Numbers}
\jednatabulkan{kravy} { |l|l |l | l | }
{
         \hline
singular &
one cow
&
ena krava
&
jedna kráva
\\
   \hline
dual (in Slovenian) &
two cows
&
dve kravi
&
dvě krávy
\\
   \hline
plural &
three cows
&
tri krave
&
tři krávy
\\


    \hline
} {Demonstrating duals on Slovenian} 

PIE had three numbers, singular, plural and dual. In PS, dual slowly disappeared, while still retaining some of its usage.

However, dual disappeared in most of Slavic languages later (including Czech and Russian), leaving only traces in the grammar. One of the languages where dual remained is Slovenian.

To illustrate, I have added a comparison of \uv{one cow}, \uv{two cows} and \uv{three cows} in Slovenian and Czech in Table~\ref{tabulka:kravy}.

In Czech, dual was retained in declensions of several words, like \uv{hands}; in Russian, the dual \uv{рукама} survives in some dialects, but is generally incorrect.\footnote{See \cite{offord1996using}, page 18.}

\subsection{Genders}
PIE had three genders, masculine, feminine and neutral. PS retained these genders, adding categories Personal and Animate.

\subsection{Tenses}
PIE had six tenses: present, future, aorist, imperfect, perfect and pluperfect. 

The tenses were somehow retained in PS; however, future, perfect and pluperfect were in many cases re-formed analytically --- creating more complex forms with auxiliary verb and either an infinitive or past participle -- for example, \uv{budu zpívat} (I will sing) in Czech, or \uv{буду петь} in Russian.

\subsection{Moods}
PIE had four moods: indicative, subjunctive, optative and imperative. In PS, imperative forms were replaced by the optatives, and subjunctive mood slowly became conditional.

\subsection{Voices}
\jednatabulkan{voices} { |l|l |l | l | }
{
         \hline

active & νιπτω & I wash (someone) & myji (někoho) \\ \hline
medium & νίπτομαι & I wash (myself) & myji se \\ \hline
passive & νίπτομαι & I am washed (by somebody) & jsem myt \\ \hline


} {Voices in Classic Greek vs. Czech} 

PIE had an active and a mediopassive voice. 

In PS, this was refined as a reflexive and a non-reflexive voice, with the addition of a passive voice.

To illustrate, I have added an example of Classic Greek that still retained the mediopassive voice, in Table~\ref{tabulka:voices}.\footnote{See for example \cite{greek1}, \cite{greek2}}

\subsection{Aspects}
\jednatabulkan{nosim} { |l|l |l |  }
{
         \hline
imperfective determinate &
нести́
&
nést
\\
   \hline
imperfective indeterminate &
носи́ть
&
nosit
\\
   \hline
perfective& 
понести 
&
ponést

\\


    \hline
} {Aspects in Czech and Russian} 

PIE had distinction between two aspects -- eventive and stative; eventive aspect being further divided into perfective and imperfective aspect.\footnote{See \cite{ringe2008proto}, page 24}

In PS, the stative aspect is degramatized\footnote{See \cite{andersen2013origin}}; however, the perfective/imperfective distinction became more important in PS than in other Indo-European languages. Also, the imperfective motion verbs were further split into determinate and non-determinate.

The determinate/non-determinate and perfective/imperfective distinction is retained in both Czech and Russian. See the Table~\ref{tabulka:nosim} and note, how hard would be to correctly translate the distinction into English.


\section{Drifting of Czech and Russian}
In about 5th century AD, the relative unity of Slavic started to break up, caused mainly by migration, which was partly caused by expansion to the north and east by the Eastern Slavs, partly by dissolution of Roman and Hun empires and the resulting vacuum in Central Europe.

One group of Slavs moved westwards, reaching approximately what is now Poland, Czech Republic and parts of Germany, creating the so-called West Slavic language group, which Czech language is a member of. The other group moved south to Balkan, creating a South Slavic language group, which I won't discuss further in this thesis. Eastern Slavs moved north and east, eventually creating an Eastern Slavic language group. Slavic was divided into those three groups at about 10th century AD.

The differences between the languages themselves are actually less syntactical and are more phonological, morphonological and morphological, with some broad-scale lexical changes. For that reason, I am not describing the linguistic differences of the languages any further, since it would be, again, not in a scope for this thesis. Again, I can recommend the book \cite{sussex2011slavic} for anybody interested in the deeper differences.



\chapter{Experiment setting}
\section{Statistical vs. rule-based translation systems}

Machine translation (MT) systems has historically used many different approaches. One way of classifying the approaches is on the axis of rule-based vs. statistical.

In general, we can re-use the descriptions, used by \cite{bojar}, which is as follows:
\begin{itemize}
\item rule based MT systems:
\begin{itemize}
\item use analysis, transfer and synthesis steps
\item rule-based machine translation systems use formal grammars
\item rule-based machine translation systems use hand-made dictionaries
\item has linguistic information hard-coded and therefore aren't language-agnostic
\end{itemize}
\item statistical MT systems
\begin{itemize}
    \item use more variants of outputs, rank them with some score, and choose the best one
    \item train internal dictionaries from big parallel data
    \item the core of the system is compact and its inner working less obvious
    \item uses statistics instead of linguistic rules and therefore are language-agnostic
\end{itemize}
\end{itemize}

However, with some of the modern system, the distinction is not as clear-cut as we would like, for the purpose of our comparison. 
For example, statistical MT systems like Moses can get significantly better with some added linguistic information; 
on the other hand, systems like TectoMT, which can sometimes be classified as more rule-based, actually have big modules more or less based on statistics.

%In this thesis, I will present several existing solutions to Czech-Russian machine translation, describe their position on the aforementioned axis, present their history and discuss their results on a testing data.
\section{Development and testing data}

To compare several machine translation systems, we need some consistent data to show the deficiencies and/or improvements of different MT systems.

In general, we want to produce a system that's general enough and performs well on unseen sentences; for that reason, we should have -- in the case of statistical MT systems -- separated test set and training/development set.\footnote{As noted in for example \cite{koehn2010statistical}, p. 274, or \cite{bishop}, pg. 6}

We have tried to do that in our experiments with our data models. However, some systems evaluated in this thesis are only black-box systems (Google Translate, PC Translator). Especially with Google Translate, we cannot rule out the possibility that they already use some parts of the test/development data.





% Ukázka použití některých konstrukcí LateXu (odkomentujte, chcete-li)
% \include{example}


\chapter*{Conclusion}
\addcontentsline{toc}{chapter}{Conclusion}
TODO


%%% Seznam použité literatury
%%% Seznam použité literatury je zpracován podle platných standardů. Povinnou citační
%%% normou pro diplomovou práci je ISO 690. Jména časopisů lze uvádět zkráceně, ale jen
%%% v kodifikované podobě. Všechny použité zdroje a prameny musí být řádně citovány.

\def\bibname{Bibliography}
%\begin{thebibliography}{99}
\addcontentsline{toc}{chapter}{\bibname}


\printbibliography



%\bibitem{lamport94}
%  {\sc Lamport,} Leslie.
%  \emph{\LaTeX: A Document Preparation System}.
%  2. vydání.
%  Massachusetts: Addison Wesley, 1994.
%  ISBN 0-201-52983-1.

%\end{thebibliography}


%%% Tabulky v diplomové práci, existují-li.
\chapwithtoc{List of Tables}

%%% Použité zkratky v diplomové práci, existují-li, včetně jejich vysvětlení.
\chapwithtoc{List of Abbreviations}

%%% Přílohy k diplomové práci, existují-li (různé dodatky jako výpisy programů,
%%% diagramy apod.). Každá příloha musí být alespoň jednou odkazována z vlastního
%%% textu práce. Přílohy se číslují.
\chapwithtoc{Attachments}

\openright
\end{document}
